\section{Open loci defined by module maps}
\label{section-loci-maps}

\noindent
The set of primes where a given module map is surjective, or an isomorphism
is sometimes open. In the case of finite projective modules we can look
at the rank of the map.

\begin{lemma}
\label{lemma-map-between-finite}
Let $R$ be a ring. Let $\varphi : M \to N$ be a map of $R$-modules
with $N$ a finite $R$-module. Then we have the equality
\begin{align*}
U & = \{\mathfrak p \subset R \mid
\varphi_{\mathfrak p} : M_{\mathfrak p} \to N_{\mathfrak p}
\text{ is surjective}\} \\
& = \{\mathfrak p \subset R \mid
\varphi \otimes \kappa(\mathfrak p) :
M \otimes \kappa(\mathfrak p) \to N \otimes \kappa(\mathfrak p)
\text{ is surjective}\}
\end{align*}
and $U$ is an open subset of $\Spec(R)$. Moreover, for any $f \in R$
such that $D(f) \subset U$ the map $M_f \to N_f$ is surjective.
\end{lemma}

\begin{proof}
The equality in the displayed formula follows from Nakayama's lemma.
Nakayama's lemma also implies that $U$ is open. See
Lemma \ref{lemma-NAK} especially part (3). If $D(f) \subset U$, then
$M_f \to N_f$ is surjective on all localizations at primes of
$R_f$, and hence it is surjective by Lemma \ref{lemma-characterize-zero-local}.
\end{proof}

\begin{lemma}
\label{lemma-map-between-finitely-presented}
Let $R$ be a ring. Let $\varphi : M \to N$ be a map of $R$-modules
with $M$ finite and $N$ finitely presented. Then
$$
U = \{\mathfrak p \subset R \mid
\varphi_{\mathfrak p} : M_{\mathfrak p} \to N_{\mathfrak p}
\text{ is an isomorphism}\}
$$
is an open subset of $\Spec(R)$.
\end{lemma}

\begin{proof}
Let $\mathfrak p \in U$. Pick a presentation
$N = R^{\oplus n}/\sum_{j = 1, \ldots, m} R k_j$.
Denote $e_i$ the image in $N$ of the $i$th basis vector of $R^{\oplus n}$.
For each $i \in \{1, \ldots, n\}$ choose an element
$m_i \in M_{\mathfrak p}$ such that $\varphi(m_i) = f_i e_i$
for some $f_i \in R$, $f_i \not \in \mathfrak p$. This is possible
as $\varphi_{\mathfrak p}$ is an isomorphism. Set $f = f_1 \ldots f_n$
and let $\psi : R_f^{\oplus n} \to M_f$ be the map which maps
the $i$th basis vector to $m_i/f_i$. Note that
$\varphi_f \circ \psi$ is the localization
at $f$ of the given map $R^{\oplus n} \to N$.
As $\varphi_{\mathfrak p}$ is an isomorphism we
see that $\psi(k_j)$ is an element of $M$ which maps to zero
in $M_{\mathfrak p}$. Hence we see that there exist
$g_j \in R$, $g_j \not \in \mathfrak p$ such that $g_j \psi(k_j) = 0$.
Setting $g = g_1 \ldots g_m$, we see that $\psi_g$ factors through
$N_{fg}$ to give a map $\chi : N_{fg} \to M_{fg}$. By construction
$\chi$ is a right inverse to $\varphi_{fg}$. It follows that
$\chi_\mathfrak p$ is an isomorphism. By Lemma \ref{lemma-map-between-finite}
there is an $h \in R$, $h \not \in \mathfrak p$ such that
$\chi_h : N_{fgh} \to M_{fgh}$ is surjective.
Hence $\varphi_{fgh}$ and $\chi_h$ are mutually inverse maps,
which implies that $D(fgh) \subset U$ as desired.
\end{proof}

\begin{lemma}
\label{lemma-finitely-presented-localization-free}
Let $R$ be a ring. Let $\mathfrak p \subset R$ be a prime.
Let $M$ be a finitely presented $R$-module. If $M_\mathfrak p$
is free, then there is an $f \in R$, $f \not \in \mathfrak p$
such that $M_f$ is a free $R_f$-module.
\end{lemma}

\begin{proof}
Choose a basis $x_1, \ldots, x_n \in M_\mathfrak p$. We can choose
an $f \in R$, $f \not \in \mathfrak p$ such that $x_i$ is the image
of some $y_i \in M_f$. After replacing $y_i$ by $f^m y_i$ for $m \gg 0$
we may assume $y_i \in M$. Namely, this replaces $x_1, \ldots, x_n$
by $f^mx_1, \ldots, f^mx_n$ which is still a basis as $f$ maps to
a unit in $R_\mathfrak p$. Hence
we obtain a homomorphism $\varphi = (y_1, \ldots, y_n) : R^{\oplus n} \to M$
of $R$-modules whose localization at $\mathfrak p$ is an isomorphism.
By Lemma \ref{lemma-map-between-finitely-presented}
we can find an $f \in R$, $f \not \in \mathfrak p$
such that $\varphi_\mathfrak q$ is an isomorphism for all primes
$\mathfrak q \subset R$ with $f \not \in \mathfrak q$.
Then it follows from Lemma \ref{lemma-characterize-zero-local}
that $\varphi_f$ is an isomorphism and the proof is complete.
\end{proof}

\begin{lemma}
\label{lemma-cokernel-flat}
Let $R$ be a ring. Let $\varphi : P_1 \to P_2$ be a map of
finite projective modules. Then
\begin{enumerate}
\item The set $U$ of primes $\mathfrak p \in \Spec(R)$ such that
$\varphi \otimes \kappa(\mathfrak p)$ is injective is open and
for any $f\in R$ such that $D(f) \subset U$ we have
\begin{enumerate}
\item $P_{1, f} \to P_{2, f}$ is injective, and
\item the module $\Coker(\varphi)_f$ is finite projective over $R_f$.
\end{enumerate}
\item The set $W$ of primes $\mathfrak p \in \Spec(R)$ such that
$\varphi \otimes \kappa(\mathfrak p)$ is surjective is open and
for any $f\in R$ such that $D(f) \subset W$ we have
\begin{enumerate}
\item $P_{1, f} \to P_{2, f}$ is surjective, and
\item the module $\Ker(\varphi)_f$ is finite projective over $R_f$.
\end{enumerate}
\item The set $V$ of primes $\mathfrak p \in \Spec(R)$ such that
$\varphi \otimes \kappa(\mathfrak p)$ is an isomorphism is open and
for any $f\in R$ such that $D(f) \subset V$ the map
$\varphi : P_{1, f} \to P_{2, f}$ is an isomorphism of modules over $R_f$.
\end{enumerate}
\end{lemma}

\begin{proof}
To prove the set $U$ is open we may work locally on
$\Spec(R)$. Thus we may replace $R$ by a suitable localization
and assume that $P_1 = R^{n_1}$ and $P_2 = R^{n_2}$, see Lemma
\ref{lemma-finite-projective}. In this case injectivity of
$\varphi \otimes \kappa(\mathfrak p)$ is equivalent to $n_1 \leq n_2$
and some $n_1 \times n_1$ minor $f$ of the matrix of $\varphi$ being
invertible in $\kappa(\mathfrak p)$. Thus $D(f) \subset U$.
This argument also shows that $P_{1, \mathfrak p} \to P_{2, \mathfrak p}$
is injective for $\mathfrak p \in U$.

\medskip\noindent
Now suppose $D(f) \subset U$. By the remark in the previous paragraph
and Lemma \ref{lemma-characterize-zero-local} we see that
$P_{1, f} \to P_{2, f}$ is injective, i.e., (1)(a) holds.
By Lemma \ref{lemma-finite-projective} to prove (1)(b)
it suffices to prove that $\Coker(\varphi)$ is finite projective
locally on $D(f)$. Thus, as we saw above, we may
assume that $P_1 = R^{n_1}$ and $P_2 = R^{n_2}$
and that some minor of the matrix of $\varphi$ is invertible in $R$.
If the minor in question corresponds to the first $n_1$
basis vectors of $R^{n_2}$, then using the last $n_2 - n_1$ basis
vectors we get a map $R^{n_2 - n_1} \to R^{n_2} \to \Coker(\varphi)$
which is easily seen to be an isomorphism.

\medskip\noindent
Openness of $W$ and (2)(a) for $D(f) \subset W$ follow from
Lemma \ref{lemma-map-between-finite}. Since $P_{2, f}$ is projective
over $R_f$ we see that $\varphi_f : P_{1, f} \to P_{2, f}$ has a section
and it follows that $\Ker(\varphi)_f$ is a direct summand of $P_{2, f}$.
Therefore $\Ker(\varphi)_f$ is finite projective. Thus (2)(b) holds as well.

\medskip\noindent
It is clear that $V = U \cap W$ is open and the other statement in (3)
follows from (1)(a) and (2)(a).
\end{proof}






